\chapter{Resultats experimentaux et valeur operationnelle}

\section{Cadre d'evaluation}
Cette section presente un premier jeu de resultats issus des artefacts du projet.
Le run analyse (\texttt{81fc268b-fe7d-49f3-86cb-07d17598686e}) a ete enregistre
le 6 mars 2026 en mode \textit{walk\_forward}. Les chiffres sont reproduits depuis
les fichiers \texttt{leaderboard\_clean.json} et \texttt{run\_details.json}.

\section{Resultats par horizon}
Le Tableau~\ref{tab:results-horizon} compare les sorties short/medium/long pour
le symbole IAM.

\begin{table}[H]
\centering
\caption{Extrait des resultats par horizon (run du 6 mars 2026)}
\label{tab:results-horizon}
\begin{tabular}{lrrrr}
\toprule
\textbf{Horizon} & \textbf{PnL} & \textbf{Total Return} & \textbf{Efficiency} & \textbf{n\_fills} \\
\midrule
Short & 14549.33 & 0.0153 & 0.1456 & 1 \\
Medium & 13776.37 & 0.1378 & 0.1379 & 1 \\
Long & 15271.22 & 0.1527 & 1.0000 & 6 \\
\bottomrule
\end{tabular}
\end{table}

Le Tableau~\ref{tab:results-horizon} montre que l'horizon long conserve la
meilleure efficience dans cet echantillon de demonstration.

\section{Lecture des resultats}
Les resultats indiquent trois points :
\begin{enumerate}
\item la plateforme produit des sorties comparables par horizon avec un format
contractuel unique ;
\item les metriques exposees permettent une lecture \textit{performance + stabilite}
plutot qu'une lecture PnL seule ;
\item la decision finale peut etre argumentee avec des traces
reproductibles (parametres, periode, scores, artefacts).
\end{enumerate}

\section{Valeur operationnelle pour le desk}
La valeur creee par la plateforme est synthetisee dans le
Tableau~\ref{tab:operational-value}.

\begin{table}[H]
\centering
\caption{Valeur operationnelle de la solution}
\label{tab:operational-value}
\begin{tabular}{p{4.5cm}p{8.8cm}}
\toprule
\textbf{Besoin desk} & \textbf{Apport de la plateforme} \\
\midrule
Comparer plusieurs hypotheses & Leaderboards et diagnostics homogenes par horizon/symbole. \\
Eviter les decisions opaques & Contrats inter-stations et parametres traces par run. \\
Limiter le sur-ajustement & Validation WFO + OOS + controles statistiques. \\
Industrialiser progressivement & Architecture modulaire compatible extensions familles/risque. \\
\bottomrule
\end{tabular}
\end{table}

Le Tableau~\ref{tab:operational-value} montre que l'outil fournit une aide a la
decision structuree, au-dela d'un simple simulateur de performance.

\section{Discussion critique}
Ces resultats restent un socle de validation technique. Pour une validation
metier complete, il faudra etendre l'analyse a un univers d'actions plus large,
integrer des periodes de stress supplementaires, et documenter un protocole
d'acceptation de production par lot de strategies.

